% Distilled blueprint for Evans, Chapter 11 (Systems of Conservation Laws).
% Formal declarations, metadata, and citation identifiers are preserved.

\section{Introduction}\label{sec:conservation-laws-introduction}

Let $\mathbf F:\mathbb R^m\to\mathbb R^m$ be smooth and let
$\mathbf g:\mathbb R\to\mathbb R^m$. The one-dimensional initial-value problem is
\[
(5)\qquad
\begin{cases}
\mathbf u_t+\mathbf F(\mathbf u)_x=0 & \text{in }\mathbb R\times(0,\infty),\\
\mathbf u=\mathbf g & \text{on }\mathbb R\times\{t=0\}.
\end{cases}
\]

\begin{example}[The $p$-system]
\label{ex:p-system-definition}
\dcref{ch11:11.1-ex-1}
\group{Conservation Laws}
The $p$-system is

$$(6) \qquad \begin{cases} u^1_t - u^2_x = 0 & \text{(compatibility condition)} \\ u^2_t - p(u^1)_x = 0 & \text{(Newton's law)}, \end{cases}$$

in $\mathbb{R}\times(0,\infty)$, where $p:\mathbb{R}\to\mathbb{R}$, with

$$(7) \qquad \mathbf{F}(z) = (-z_2, -p(z_1))$$

for $z=(z_1,z_2)$. Under $u^1=u_x$ and $u^2=u_t$, it is equivalent to

$$(8) \qquad u_{tt} - (p(u_x))_x = 0 \quad \text{in } \mathbb{R}\times(0,\infty).$$

\end{example}

\begin{example}[Euler's equations for compressible gas flow]
\label{ex:euler-equations-1d-system}
\dcref{ch11:11.1-ex-2}
\group{Conservation Laws}
\textit{Euler's equations} for compressible gas flow in one dimension are

$$(9) \qquad \begin{cases} \rho_t + (\rho v)_x = 0 & \text{(conservation of mass)} \\ (\rho v)_t + (\rho v^2+p)_x = 0 & \text{(conservation of momentum)} \\ (\rho E)_t + (\rho E v + pv)_x = 0 & \text{(conservation of energy)} \end{cases}$$

in $\mathbb{R}\times(0,\infty)$. Here $\rho$ is the \textit{mass density}, $v$ the \textit{velocity}, and $E$ the \textit{energy density} per unit mass. We assume

$$E = e + \frac{v^2}{2},$$

where $e$ is the \textit{internal energy} per unit mass and the term $\frac{v^2}{2}$ corresponds to the kinetic energy per unit mass. The letter $p$ in (9) denotes the \textit{pressure}. We assume $p$ is a known function

$$(10) \qquad p = p(\rho,e)$$

of $\rho$ and $e$; formula (10) is a \textit{constitutive relation}. Writing $\mathbf{u}=(u^1,u^2,u^3)=(\rho,\rho v,\rho E)$, we check (9) is a system of conservation laws of the requisite form
$$\mathbf{u}_t + \mathbf{F}(\mathbf{u})_x = 0 \quad \text{in } \mathbb{R}\times(0,\infty)$$
for $\mathbf{F}=(F^1,F^2,F^3)$,

$$(11) \qquad \begin{cases} F^1(z) = z_2 \\ F^2(z) = \frac{(z_2)^2}{z_1} + p\left(z_1, \frac{z_3}{z_1} - \frac{1}{2}\left(\frac{z_2}{z_1}\right)^2\right) \\ F^3(z) = \frac{z_2 z_3}{z_1} + p\left(z_1, \frac{z_3}{z_1} - \frac{1}{2}\left(\frac{z_2}{z_1}\right)^2\right)\frac{z_2}{z_1}, \end{cases}$$

where $z=(z_1,z_2,z_3)$ and $z_1 > 0$.
\end{example}

\begin{example}[The shallow water equations]
\label{ex:shallow-water-equations}
\dcref{ch11:11.1-ex-3}
\group{Conservation Laws}
The one-dimensional \textit{shallow water equations} are
$$\begin{cases} \phi_t + (v\phi)_x = 0 & \text{(conservation of mass)} \\ v_t + (v^2/2+\phi)_x = 0 & \text{(conservation of momentum)} \end{cases}$$
in $\mathbb{R}\times[0,\infty)$. In these equations $v$ is the horizontal velocity, $\phi$ denotes $gh$, where $h$ is the height and $g$ the gravitational constant, and
$$\mathbf{F}(z) = \left(z_1z_2, \frac{(z_2)^2}{2} + z_1\right).$$
\end{example}

\subsection{Integral solutions}\label{sec:integral-solutions-systems}

For every smooth, compactly supported
\[
(12)\qquad \mathbf v:\mathbb R\times[0,\infty)\to\mathbb R^m,
\]
the weak formulation of (5) is
\[
(13)\qquad
\int_0^\infty\!\int_{-\infty}^{\infty}
\bigl(\mathbf u\cdot\mathbf v_t+\mathbf F(\mathbf u)\cdot\mathbf v_x\bigr)\,dx\,dt
+\int_{-\infty}^{\infty}\mathbf g\cdot\mathbf v\big|_{t=0}\,dx=0.
\]

\begin{definition}[Integral solution of a system of conservation laws]
\label{def:integral-solution-system}
\uses{def:integral-solution-conservation-law}
\dcref{ch11:11.1-def-1}
\group{Conservation Laws}

An \textit{integral solution} of (5) is a function
$\mathbf u\in L^\infty(\mathbb R\times(0,\infty);\mathbb R^m)$ satisfying (13) for every test function (12).
\end{definition}

\begin{proposition}[Rankine--Hugoniot jump condition]
\label{prop:rankine-hugoniot-condition-system}
\uses{def:integral-solution-system}
\dcref{ch11:11.1-prop-1}
\group{Conservation Laws}

Suppose $\mathbf u$ is an integral solution, smooth on either side of the smooth curve
$C=\{(x,t)\mid x=s(t)\}$. Then, along $C$,
$$(20) \qquad [\![\mathbf{F}(\mathbf{u})]\!] = \sigma[\![\mathbf{u}]\!],$$
where
$$\begin{cases}
[\![\mathbf{u}]\!] = \mathbf{u}_l - \mathbf{u}_r = \text{jump in } \mathbf{u} \text{ across the curve } C \\
[\![\mathbf{F}(\mathbf{u})]\!] = \mathbf{F}(\mathbf{u}_l) - \mathbf{F}(\mathbf{u}_r) = \text{jump in } \mathbf{F}(\mathbf{u}) \\
\sigma = \dot{s} = \text{speed of the curve } C,
\end{cases}$$
the subscript ``$l$'' denoting the limit from the left and ``$r$'' the limit from the right.
\end{proposition}

\begin{proof}
\sketch Apply (13) to test functions supported near $C$ and integrate by parts separately on its two sides. The interior terms vanish by the classical equation, while the boundary terms give
$n_t[\![\mathbf u]\!]+n_x[\![\mathbf F(\mathbf u)]\!]=0$ for a normal $(n_x,n_t)$ to $C$. Since a normal to $x=s(t)$ is proportional to $(1,-\dot s)$, this is (20).
\end{proof}

\subsection{Traveling waves, hyperbolic systems}\label{sec:traveling-waves-hyperbolic-systems}

Consider the quasilinear system
\[
(21)\qquad \mathbf u_t+\mathbf B(\mathbf u)\mathbf u_x=0,
\]
where $\mathbf B:\mathbb R^m\to\mathbb R^{m\times m}$ is smooth; a conservation law has $\mathbf B=D\mathbf F$. A traveling wave $\mathbf u(x,t)=\mathbf v(x-\sigma t)$ satisfies
\[
(23)\qquad (\mathbf B(\mathbf v)-\sigma I)\mathbf v'=0.
\]

\begin{definition}[Strict hyperbolicity]
\label{def:strictly-hyperbolic-system}
\dcref{ch11:11.1-def-2}
\group{Conservation Laws}
If for each $z\in\mathbb{R}^m$ the eigenvalues of $\mathbf{B}(z)$ are real and distinct, we call the system (21) \textit{strictly hyperbolic}.
\end{definition}

\begin{definition}[Eigenvalues and eigenvectors of a hyperbolic system]
\label{def:eigenvalue-eigenvector-notation-hyperbolic}
\uses{def:strictly-hyperbolic-system}
\dcref{ch11:11.1-def-3}
\group{Conservation Laws}

(i) Order the eigenvalues as
$$(24) \qquad \lambda_1(z) < \lambda_2(z) < \cdots < \lambda_m(z) \quad (z\in\mathbb{R}^m)$$
for each $z\in\mathbb R^m$.

(ii) For each $k=1,\ldots,m$, choose a nonzero right eigenvector
$$\mathbf{r}_k(z)$$
satisfying
$$(25) \qquad \mathbf{B}(z)\mathbf{r}_k(z) = \lambda_k(z)\mathbf{r}_k(z) \quad (k=1,\ldots,m,\ z\in\mathbb{R}^m).$$
Since we are always assuming the strict hyperbolicity condition, the vectors $\{\mathbf{r}_k(z)\}_{k=1}^m$ span $\mathbb{R}^m$ for each $z\in\mathbb{R}^m$.

(iii) For each $k=1,\ldots,m$, choose a nonzero eigenvector
$$\mathbf{l}_k(z)$$
of $\mathbf B(z)^T$ corresponding to $\lambda_k(z)$. Thus
$$(26) \qquad \mathbf{B}(z)^T\mathbf{l}_k(z) = \lambda_k(z)\mathbf{l}_k(z) \quad (k=1,\ldots,m,\ z\in\mathbb{R}^m).$$
This equality is usually written
$$(27) \qquad \mathbf{l}_k(z)\mathbf{B}(z) = \lambda_k(z)\mathbf{l}_k(z) \quad (k=1,\ldots,m,\ z\in\mathbb{R}^m).$$
Thus $\{\mathbf{l}_k(z)\}_{k=1}^m$ are the \textit{left eigenvectors} and $\{\mathbf{r}_k(z)\}_{k=1}^m$ the \textit{right eigenvectors}.
\end{definition}

\begin{remark}[Orthogonality of left and right eigenvectors]
\label{rem:eigenvector-orthogonality}
\uses{def:eigenvalue-eigenvector-notation-hyperbolic}
\dcref{ch11:11.1-rem-1}
\group{Conservation Laws}

$$(28) \qquad \mathbf{l}_l(z)\cdot\mathbf{r}_k(z) = 0 \quad \text{if } k\neq l \quad (z\in\mathbb{R}^m).$$
Indeed, (25) and (26) give
$$\lambda_k(z)(\mathbf{l}(z)\cdot\mathbf{r}_k(z)) = \mathbf{l}(z)\cdot(\mathbf{B}(z)\mathbf{r}_k(z)) = (\mathbf{B}(z)^T\mathbf{l}(z))\cdot\mathbf{r}_k(z) = \lambda_l(z)(\mathbf{l}(z)\cdot\mathbf{r}_k(z));$$
strict hyperbolicity implies (28).
\end{remark}

\begin{theorem}[Invariance of hyperbolicity under change of coordinates]
\label{thm:invariance-hyperbolicity-coordinate-change}
\uses{def:strictly-hyperbolic-system, def:eigenvalue-eigenvector-notation-hyperbolic}
\dcref{ch11:11.1-thm-1}
\group{Conservation Laws}

Let $u$ be a smooth solution of the strictly hyperbolic system (21). Assume also $\Phi:\mathbb{R}^m\to\mathbb{R}^m$ is a smooth diffeomorphism, with inverse $\Psi$. Then
$$(29) \qquad \tilde{u} := \Phi(u)$$
solves the strictly hyperbolic system
$$(30) \qquad \tilde u_t + \tilde B(\tilde u)\tilde u_x = 0 \quad \text{in } \mathbb{R}\times(0,\infty),$$
for
$$(31) \qquad \tilde B(z) := D\Phi(\Psi(z))B(\Psi(z))D\Psi(z) \quad (z\in\mathbb{R}^m).$$
\end{theorem}

\begin{proof}
\sketch The chain rule gives $\widetilde u_t=D\Phi(u)u_t$ and $\widetilde u_x=D\Phi(u)u_x$. Multiplying (21) by $D\Phi(u)$ and using $D\Psi(\widetilde u)D\Phi(u)=I$ yields (30)--(31). The matrices $\widetilde B(\Phi(u))$ and $B(u)$ are similar, so they have the same real, distinct eigenvalues; moreover $D\Phi(u)r_k(u)$ is a corresponding right eigenvector.
\end{proof}

\begin{theorem}[Dependence of eigenvalues and eigenvectors on parameters]
\label{thm:smooth-dependence-eigenvalues-eigenvectors}
\uses{def:strictly-hyperbolic-system, def:eigenvalue-eigenvector-notation-hyperbolic}
\dcref{ch11:11.1-thm-2}
\group{Conservation Laws}

Assume the matrix function $\mathbf{B}$ is smooth, strictly hyperbolic. Then

(i) the eigenvalues $\lambda_k(z)$ depend smoothly on $z\in\mathbb{R}^m$ $(k=1,\ldots,m)$.

(ii) Furthermore, we can select the right eigenvectors $\mathbf{r}_k(z)$ and left eigenvectors $\mathbf{l}_k(z)$ to depend smoothly on $z\in\mathbb{R}^m$ and satisfy the normalization
$$|\mathbf{r}_k(z)|,\ |\mathbf{l}_k(z)| = 1 \quad (k=1,\ldots,m).$$
\end{theorem}

\begin{proof}
\sketch Fix $z_0$ and a simple eigenpair $(\lambda_k(z_0),r_k(z_0))$. Apply the implicit function theorem to $B(z)r-\lambda r=0$ together with a local normalization fixing the scale of $r$. Simplicity makes the linearized eigenpair equations invertible, producing smooth local eigenvalues and normalized eigenvectors. Ordering the distinct eigenvalues patches the local eigenvalue branches. The one-dimensional eigenspaces allow the normalized eigenvectors to be continued with fixed orientation along radial paths, and uniqueness patches these choices globally. Apply the same argument to $B^T$ for the left eigenvectors.
\end{proof}

\begin{remark}[Global orientation of eigenspaces]
\label{rem:global-orientation-eigenspaces}
\uses{thm:smooth-dependence-eigenvalues-eigenvectors}
\dcref{ch11:11.1-rem-2}
\group{Conservation Laws}

The globally smooth normalized eigenvectors orient the one-dimensional eigenspaces. This conclusion uses strict hyperbolicity and can fail for higher-dimensional eigenspaces.
\end{remark}

\begin{example}[Strict hyperbolicity of the $p$-system]
\label{ex:p-system-strict-hyperbolicity}
\uses{ex:p-system-definition, def:strictly-hyperbolic-system}
\dcref{ch11:11.1-ex-4}
\group{Conservation Laws}

For the $p$-system (6), we have
$$\mathbf{u}_t + \mathbf{B}(\mathbf{u})\mathbf{u}_x = 0,$$
for
$$\mathbf{B}(z) = DF(z) = \begin{pmatrix} 0 & -1 \\ -p'(z_1) & 0 \end{pmatrix}.$$
The eigenvalues are $\lambda_1=-\sigma$, $\lambda_2=\sigma$, for $\sigma:=p'(z_1)^{1/2}$. These are real and distinct provided we hereafter suppose the strict hyperbolicity condition
$$(41) \qquad p'>0.$$
For (8), condition (41) says that the stress $p(u_x)$ is strictly increasing in the strain $u_x$.
\end{example}

\begin{example}[Strict hyperbolicity of Euler's equations]
\label{ex:euler-equations-strict-hyperbolicity}
\uses{ex:euler-equations-1d-system, thm:invariance-hyperbolicity-coordinate-change, def:strictly-hyperbolic-system}
\dcref{ch11:11.1-ex-5}
\group{Conservation Laws}

Euler's equations (9) comprise a strictly hyperbolic system provided we assume $p>0$ and
$$(42) \qquad \frac{\partial p}{\partial\rho}>0, \quad \frac{\partial p}{\partial e}>0,$$
where $p=p(\rho,e)$ is the constitutive relation between the mass density, the internal energy density and the pressure. This assertion is however difficult to verify directly, as the flux function $\mathbf{F}$ defined by (11) is complicated.

In the variables density $\rho$, velocity $v$, and internal energy $e$, Euler's equations are equivalently
$$(43) \qquad \begin{cases} \rho_t + v\rho_x + \rho v_x = 0 \\ v_t + vv_x + \frac{1}{\rho}p_x = 0 \\ e_t + ve_x + \frac{v}{\rho}v_x = 0, \end{cases}$$
provided $\rho>0$. These equations are not in conservation form. Setting now $\mathbf{u}=(u^1,u^2,u^3)=(\rho,v,e)$, we rewrite (43) as
$$(44) \qquad \mathbf{u}_t + \mathbf{B}(\mathbf{u})\mathbf{u}_x = 0 \quad \text{in } \mathbb{R}\times(0,\infty),$$
for
$$(45) \qquad \mathbf{B}(z) = z_2 I + \tilde{\mathbf{B}}(z),$$
where
$$\tilde{\mathbf{B}}(z) := \begin{pmatrix} 0 & z_1 & 0 \\ \frac{1}{z_1}\frac{\partial p}{\partial\rho}(z_1,z_3) & 0 & \frac{1}{z_1}\frac{\partial p}{\partial\varepsilon}(z_1,z_3) \\ 0 & \frac{1}{z_1}p(z_1,z_3) & 0 \end{pmatrix}.$$
The characteristic polynomial of $\mathbf{B}$ is $-\lambda(\lambda^2-\sigma^2)$, for $\sigma^2 = \frac{p}{\rho^2}\frac{\partial p}{\partial\varepsilon} + \frac{\partial p}{\partial\rho}$. By (45), in physical notation the eigenvalues of $\mathbf B$ are
$$(46) \qquad \lambda_1 = v-\sigma, \quad \lambda_2 = v, \quad \lambda_3 = v+\sigma,$$
where
$$\sigma := \left(\frac{p}{\rho^2}\frac{\partial p}{\partial\varepsilon} + \frac{\partial p}{\partial\rho}\right)^{1/2} > 0$$
is the local \textit{sound speed}. Thus (42) makes (44) strictly hyperbolic, and \cref{thm:invariance-hyperbolicity-coordinate-change} gives the same conclusion and eigenvalues for (9).
\end{example}

\section{Riemann's Problem}\label{sec:riemanns-problem}

Riemann's problem is
\[
(1)\qquad \mathbf u_t+\mathbf F(\mathbf u)_x=0
\quad\text{in }\mathbb R\times(0,\infty),
\qquad
(2)\qquad
\mathbf g(x)=
\begin{cases}
\mathbf u_l & x<0,\\
\mathbf u_r & x>0.
\end{cases}
\]

\subsection{Simple waves}\label{sec:simple-waves}

A $k$-simple wave has the form
\[
(3)\qquad \mathbf u(x,t)=\mathbf v(w(x,t)),
\]
where
\[
(5)\qquad w_t+\lambda_k(\mathbf v(w))w_x=0,
\qquad
(6)\qquad \mathbf v'(s)=\mathbf r_k(\mathbf v(s)).
\]
Equivalently, $w$ solves the scalar conservation law
\[
(7)\qquad w_t+F_k(w)_x=0,
\qquad
(8)\qquad F_k(s)=\int_0^s\lambda_k(\mathbf v(t))\,dt.
\]

\begin{definition}[Rarefaction curve]
\label{def:rarefaction-curve}
\uses{def:eigenvalue-eigenvector-notation-hyperbolic}
\dcref{ch11:11.2-def-1}
\group{Conservation Laws}

Given a fixed state $z_0\in\mathbb{R}^m$, we define the $k^{\text{th}}$-\textit{rarefaction curve}
$$R_k(z_0)$$
to be the path in $\mathbb{R}^m$ of the solution of the ODE (6) which passes through $z_0$.
\end{definition}

For a genuinely nonlinear pair, set
\[
R_k^+(z_0)=\{z\in R_k(z_0)\mid\lambda_k(z)>\lambda_k(z_0)\},
\qquad
R_k^-(z_0)=\{z\in R_k(z_0)\mid\lambda_k(z)<\lambda_k(z_0)\}.
\]
Then $R_k(z_0)=R_k^+(z_0)\cup\{z_0\}\cup R_k^-(z_0)$.

\begin{definition}[Genuine nonlinearity and linear degeneracy]
\label{def:genuinely-nonlinear-linearly-degenerate}
\uses{def:eigenvalue-eigenvector-notation-hyperbolic}
\dcref{ch11:11.2-def-2}
\group{Conservation Laws}

(i) The pair $(\lambda_k(z), r_k(z))$ is called \textit{genuinely nonlinear} provided
$$(11) \qquad D\lambda_k(z)\cdot r_k(z) \neq 0 \quad \text{for all } z\in\mathbb{R}^m.$$

(ii) We say $(\lambda_k(z), r_k(z))$ is \textit{linearly degenerate} if
$$(12) \qquad D\lambda_k(z)\cdot r_k(z) = 0 \quad \text{for all } z\in\mathbb{R}^m.$$
\end{definition}

\subsection{Rarefaction waves}\label{sec:rarefaction-waves}

\begin{theorem}[Existence of $k$-rarefaction waves]
\label{thm:existence-k-rarefaction-waves}
\uses{def:genuinely-nonlinear-linearly-degenerate, def:rarefaction-curve, thm:riemann-problem-solution}
\dcref{ch11:11.2-thm-1}
\group{Conservation Laws}

Suppose that for some $k\in\{1,\ldots,m\}$,

(i) the pair $(\lambda_k, r_k)$ is genuinely nonlinear, and

(ii) $u_r \in R_k^+(u_l)$.

Then there exists a continuous integral solution $\mathbf{u}$ of Riemann's problem (1), (2), which is a $k$-simple wave constant along lines through the origin.
\end{theorem}

\begin{proof}
\sketch Parameterize the rarefaction curve by $u_l=\mathbf v(w_l)$ and $u_r=\mathbf v(w_r)$. Since
$F_k''(s)=D\lambda_k(\mathbf v(s))\cdot r_k(\mathbf v(s))$, genuine nonlinearity and $u_r\in R_k^+(u_l)$ make $F_k$ strictly convex when $w_l<w_r$ and strictly concave when $w_l>w_r$. The scalar Riemann problem (7) with states $w_l,w_r$ therefore has a continuous centered rarefaction solution $w(x/t)$. Composing with $\mathbf v$ gives the required continuous integral solution of (1)--(2).
\end{proof}

\begin{remark}[Centered $k$-rarefaction wave]
\label{rem:k-rarefaction-wave-terminology}
\uses{thm:existence-k-rarefaction-waves}
\dcref{ch11:11.2-rem-1}
\group{Conservation Laws}

We call $\mathbf{u}$ a (centered) $k$-\textit{rarefaction wave}.
\end{remark}

\subsection{Shock waves, contact discontinuities}\label{sec:shock-waves-contact-discontinuities}

\subsubsection*{a. The shock set}

\begin{definition}[Shock set]
\label{def:shock-set}
\uses{prop:rankine-hugoniot-condition-system}
\dcref{ch11:11.2-def-3}
\group{Conservation Laws}

Given a fixed state $z_0\in\mathbb{R}^m$, we define the \textit{shock set}
$$S(z_0) := \{z\in\mathbb{R}^m \mid \mathbf{F}(z) - \mathbf{F}(z_0) = \sigma(z-z_0) \text{ for a constant } \sigma = \sigma(z,z_0)\}.$$
\end{definition}

\begin{theorem}[Structure of the shock set]
\label{thm:structure-of-shock-set}
\uses{def:shock-set, thm:smooth-dependence-eigenvalues-eigenvectors, def:eigenvalue-eigenvector-notation-hyperbolic}
\dcref{ch11:11.2-thm-2}
\group{Conservation Laws}

Fix $z_0\in\mathbb{R}^m$. In some neighborhood of $z_0$, $S(z_0)$ is the union of $m$ smooth curves $S_k(z_0)$ $(k=1,\ldots,m)$ such that:

(i) The curve $S_k(z_0)$ passes through $z_0$, with tangent $\mathbf{r}_k(z_0)$.

(ii) $\displaystyle\lim_{\substack{z\to z_0 \\ z\in S_k(z_0)}} \sigma(z,z_0) = \lambda_k(z_0)$.

(iii) $\sigma(z,z_0) = \dfrac{\lambda_k(z)+\lambda_k(z_0)}{2} + O(|z-z_0|^2)$, as $z\to z_0$ with $z\in S_k(z_0)$.
\end{theorem}

\begin{proof}
\sketch Define the averaged Jacobian
$\overline B(z)=\int_0^1DF(z_0+t(z-z_0))\,dt$. The Rankine--Hugoniot relation is equivalent to $(\overline B(z)-\sigma I)(z-z_0)=0$. Smooth eigenpairs of $\overline B$ and the implicit function theorem applied to the $m-1$ equations
$\overline l_j(z)\cdot(z-z_0)=0$ for $j\ne k$ produce a smooth curve $S_k(z_0)$ tangent to $r_k(z_0)$, with $\sigma\to\lambda_k(z_0)$. Differentiate the Rankine--Hugoniot identity twice along $S_k$ and differentiate
$DF(\psi_k(t))r_k(\psi_k(t))=\lambda_k(\psi_k(t))r_k(\psi_k(t))$ along $R_k$. Pairing the resulting identities with $l_k(z_0)$ gives $2\sigma'(0)=\lambda_k'(0)$, which yields (iii).
\end{proof}

\begin{theorem}[Linear degeneracy]
\label{thm:linear-degeneracy-shock-rarefaction-coincide}
\uses{thm:structure-of-shock-set, def:rarefaction-curve, def:genuinely-nonlinear-linearly-degenerate}
\dcref{ch11:11.2-thm-3}
\group{Conservation Laws}

Suppose for some $k\in\{1,\ldots,m\}$ that the pair $(\lambda_k, r_k)$ is linearly degenerate. Then for each $z_0\in\mathbb{R}^m$,

(i) $R_k(z_0) = S_k(z_0)$

and

(ii) $\sigma(z,z_0) = \lambda_k(z) = \lambda_k(z_0)$ for all $z\in S_k(z_0)$.
\end{theorem}

\begin{proof}
\sketch Let $\mathbf v'(s)=r_k(\mathbf v(s))$ and $\mathbf v(0)=z_0$. Linear degeneracy makes $\lambda_k(\mathbf v(s))=\lambda_k(z_0)$. Hence
\[
\mathbf F(\mathbf v(s))-\mathbf F(z_0)
=\int_0^sDF(\mathbf v(t))r_k(\mathbf v(t))\,dt
=\lambda_k(z_0)(\mathbf v(s)-z_0).
\]
Thus the rarefaction curve lies in the $k$-shock curve with speed $\lambda_k(z_0)$; local uniqueness of the smooth curves gives equality and the stated speed identities.
\end{proof}

\subsubsection*{b. Contact discontinuities, shock waves}

For $u_r\in S_k(u_l)$, the Rankine--Hugoniot solution is
\[
\mathbf u(x,t)=
\begin{cases}
\mathbf u_l & x<\sigma t,\\
\mathbf u_r & x>\sigma t,
\end{cases}
\qquad \sigma=\sigma(u_r,u_l).
\]
If the $k$-field is linearly degenerate, then
$\sigma=\lambda_k(u_l)=\lambda_k(u_r)$ and the discontinuity is a $k$-contact discontinuity.

\begin{definition}[Admissible pair, Lax entropy condition]
\label{def:admissible-pair-lax-entropy-condition}
\uses{def:shock-set, thm:structure-of-shock-set}
\dcref{ch11:11.2-def-4}
\group{Conservation Laws}

Assume the pair $(\lambda_k, r_k)$ is genuinely nonlinear at $u_l$. We say that the pair $(u_l, u_r)$ is \textit{admissible} provided
$$u_r \in S_k(u_l)$$
and
$$(42) \qquad \lambda_k(u_r) < \sigma(u_r,u_l) < \lambda_k(u_l).$$
Condition (42) is the \textit{Lax entropy condition}; the displayed discontinuous solution is then a \textit{$k$-shock wave}.
\end{definition}

\begin{definition}[$S_k^+$ and $S_k^-$]
\label{def:shock-set-plus-minus}
\uses{def:shock-set, def:admissible-pair-lax-entropy-condition}
\dcref{ch11:11.2-def-5}
\group{Conservation Laws}

If the pair $(\lambda_k, r_k)$ is genuinely nonlinear, we write
$$S_k^+(z_0) := \{z\in S_k(z_0) \mid \lambda_k(z_0) < \sigma(z,z_0) < \lambda_k(z)\}$$
and
$$S_k^-(z_0) := \{z\in S_k(z_0) \mid \lambda_k(z) < \sigma(z,z_0) < \lambda_k(z_0)\}.$$
Then
$$S_k(z_0) = S_k^+(z_0) \cup \{z_0\} \cup S_k^-(z_0)$$
near $z_0$. Note then that the pair $(u_l, u_r)$ is admissible if and only if
$$u_r \in S_k^-(u_l).$$
\end{definition}

\subsection{Local solution of Riemann's problem}\label{sec:local-solution-riemanns-problem}

\begin{definition}[The curve $T_k$]
\label{def:t-curve}
\uses{def:rarefaction-curve, def:shock-set-plus-minus, def:genuinely-nonlinear-linearly-degenerate}
\dcref{ch11:11.2-def-6}
\group{Conservation Laws}

(i) If the pair $(\lambda_k, r_k)$ is genuinely nonlinear, write
$$T_k(z_0) := R_k^+(z_0) \cup \{z_0\} \cup S_k^-(z_0).$$

(ii) If the pair $(\lambda_k, r_k)$ is linearly degenerate, we set
$$T_k(z_0) := R_k(z_0) = S_k(z_0).$$
\end{definition}

Each $T_k(z_0)$ is $C^1$ near $z_0$ and tangent there to $r_k(z_0)$.

\begin{theorem}[Local solution of Riemann's problem]
\label{thm:local-solution-riemann-problem}
\uses{def:t-curve, thm:existence-k-rarefaction-waves, def:admissible-pair-lax-entropy-condition, def:eigenvalue-eigenvector-notation-hyperbolic}
\dcref{ch11:11.2-thm-4}
\group{Conservation Laws}

Assume for each $k=1,\ldots,m$ that the pair $(\lambda_k, r_k)$ is either genuinely nonlinear or else linearly degenerate. Suppose further the left state $u_l$ is given. Then for each right state $u_r$ sufficiently close to $u_l$ there exists an integral solution $\mathbf{u}$ of Riemann's problem, which is constant on lines through the origin.
\end{theorem}

\begin{proof}
\sketch Parameterize each $T_k$ by signed arc length. Starting from $z_0=u_l$, move successively by parameters $t_1,\ldots,t_m$ along $T_1,\ldots,T_m$, and let $\Phi(t)$ be the final state. Then $\Phi(0)=u_l$ and the columns of $D\Phi(0)$ are $r_1(u_l),\ldots,r_m(u_l)$. Strict hyperbolicity makes these vectors a basis, so the inverse function theorem supplies parameters connecting $u_l$ to every nearby $u_r$. The corresponding rarefaction waves, admissible shocks, and contact discontinuities have strictly ordered nearby characteristic speeds and therefore concatenate into an integral solution depending only on $x/t$.
\end{proof}

\section{Systems of Two Conservation Laws}\label{sec:systems-two-conservation-laws}

Let $\mathbf F=(F^1,F^2)$, $\mathbf g=(g^1,g^2)$, and $\mathbf u=(u^1,u^2)$. The initial-value problem is
\[
(1)\qquad
\begin{cases}
u^1_t+F^1(u^1,u^2)_x=0 & \text{in }\mathbb R\times(0,\infty),\\
u^2_t+F^2(u^1,u^2)_x=0,\\
u^1=g^1,\quad u^2=g^2 & \text{on }\mathbb R\times\{t=0\}.
\end{cases}
\]

\subsection{Riemann invariants}\label{sec:riemann-invariants}

\begin{definition}[$i$-th Riemann invariant]
\label{def:riemann-invariant}
\uses{def:eigenvalue-eigenvector-notation-hyperbolic}
\dcref{ch11:11.3-def-1}
\group{Conservation Laws}

We say
$$w^i:\mathbb{R}^2\to\mathbb{R}$$
is an $i^{\text{th}}$-\textit{Riemann invariant} provided
$$(2) \qquad Dw^i(z) \text{ is parallel to } l_j(z) \quad (z\in\mathbb{R}^2, i\neq j).$$
\end{definition}

In two dimensions, (2) is equivalent to
\[
(2')\qquad Dw^i(z)\cdot r_i(z)=0,
\]
or equivalently
\[
(3)\qquad w^i\text{ is constant on each }R_i.
\]

\begin{remark}[Riemann invariants for $m>2$]
\label{rem:riemann-invariants-mgt2-remark}
\uses{def:riemann-invariant}
\dcref{ch11:11.3-rem-1}
\group{Conservation Laws}

For $m>2$, Riemann invariants need not exist.
\end{remark}

Assume
\[
(4)\qquad \mathbf w(z)=(w^1(z),w^2(z))
\]
is a local coordinate map with inverse $\mathbf z(w)$, let $\mathbf u$ be a smooth solution of (1), and set
\[
(5)\qquad \mathbf v(x,t)=\mathbf w(\mathbf u(x,t)).
\]

\begin{theorem}[Conservation laws and Riemann invariants]
\label{thm:conservation-laws-riemann-invariants}
\uses{def:riemann-invariant}
\dcref{ch11:11.3-thm-1}
\group{Conservation Laws}

The functions $v^1, v^2$ solve the system
$$(6) \qquad \begin{cases} v^1_t + \lambda_2(\mathbf{u})v^1_x = 0 \\ v^2_t + \lambda_1(\mathbf{u})v^2_x = 0 \end{cases} \quad \text{in } \mathbb{R}\times(0,\infty).$$
Thus the system is diagonal in the Riemann coordinates.
\end{theorem}

\begin{proof}
\sketch For $i\ne j$, the chain rule and (1) give
\[
v^i_t+\lambda_j(\mathbf u)v^i_x
=Dw^i(\mathbf u)\bigl(-D\mathbf F(\mathbf u)+\lambda_j(\mathbf u)I\bigr)\mathbf u_x.
\]
By definition $Dw^i$ is parallel to the left eigenvector $l_j$, so the last expression vanishes.
\end{proof}

\begin{remark}[Riemann invariants along characteristics]
\label{rem:riemann-invariant-characteristic-curves}
\uses{thm:conservation-laws-riemann-invariants, def:genuinely-nonlinear-linearly-degenerate}
\dcref{ch11:11.3-rem-2}
\group{Conservation Laws}

(i) Along the characteristics
$$(7) \qquad \dot x_i(s) = \lambda_j(\mathbf{u}(x_i(s),s)) \quad (s\geq 0)$$
for $i=1,2$, $j\neq i$, equation (6) gives
$$(8) \qquad v^i \text{ is constant along the curve } (x_i(s),s) \quad (s\geq 0)$$
for $i=1,2$.

(ii) Genuine nonlinearity is
$$(9) \qquad D\lambda_i(z)\cdot\mathbf{r}_i(z) \neq 0 \quad (z\in\mathbb{R}^2, i=1,2).$$
Since we can also regard $\lambda_i$ as a function of $w=(w_1,w_2)$, we can rewrite (9) to read
$$(10) \qquad \frac{\partial\lambda_i}{\partial w_j} \neq 0 \quad (w\in\mathbb{R}^2, i\neq j).$$
Indeed, failure of (10) gives
$$(11) \qquad 0 = \frac{\partial\lambda_i}{\partial w_j} = \sum_{k=1}^2 \frac{\partial\lambda_i}{\partial z_k}\frac{\partial z^k}{\partial w_j}.$$
The inverse-coordinate identities imply that (11) is equivalent to $D\lambda_i$ being parallel to $Dw^i$, which is orthogonal to $r_i$. Hence (9) and (10) are equivalent.
\end{remark}

\begin{example}[Barotropic compressible gas dynamics]
\label{ex:barotropic-gas-dynamics-riemann-invariants}
\uses{def:riemann-invariant, ex:euler-equations-1d-system, rem:riemann-invariant-characteristic-curves}
\dcref{ch11:11.3-ex-1}
\group{Conservation Laws}

For barotropic compressible gas dynamics, the relevant equations are
$$(12) \qquad \begin{cases} \rho_t + (\rho v)_x = 0 & \text{(conservation of mass)} \\ (\rho v)_t + (\rho v^2+p)_x = 0 & \text{(conservation of momentum)}, \end{cases}$$
with the barotropic equation of state
$$(13) \qquad p = p(\rho)$$
for some smooth function $p:\mathbb{R}\to\mathbb{R}$. Formula (13) is called a \textit{barotropic equation of state}. We assume the strict hyperbolicity condition
$$(14) \qquad p' > 0.$$
Setting $\mathbf{u}=(u^1,u^2)=(\rho,\rho v)$, we can rewrite (12), (13) to read
$$\mathbf{u}_t + \mathbf{F}(\mathbf{u})_x = 0,$$
for
$$\mathbf{F} = (F^1,F^2) = (z_2, (z_2)^2/z_1 + p(z_1))$$
and $z=(z_1,z_2)$, provided $z_1>0$. Then
$$D\mathbf{F} = \begin{pmatrix} 0 & 1 \\ -(z_2/z_1)^2 + p'(z_1) & 2z_2/z_1 \end{pmatrix}.$$
Consequently
$$(15) \qquad \lambda_1 = \frac{z_2}{z_1} - p'(z_1)^{1/2}, \quad \lambda_2 = \frac{z_2}{z_1} + p'(z_1)^{1/2}.$$
In physical notation:
$$(16) \qquad \lambda_1 = v-\sigma, \quad \lambda_2 = v+\sigma,$$
for the sound speed
$$(17) \qquad \sigma := p'(\rho)^{1/2}.$$
The two characteristic families are
$$(18) \qquad \dot x_1(t) = v(x_1(t),t) + \sigma(x_1(t),t),$$
$$(19) \qquad \dot x_2(t) = v(x_2(t),t) - \sigma(x_2(t),t),$$
where $\sigma(x,t) := p'(\rho(x,t))^{\frac12}$, $t\geq 0$. We know from (8) that the Riemann invariant $w^1=w^1(\mathbf{u})$ is constant along the trajectories of (18) and $w^2=w^2(\mathbf{u})$ is constant along trajectories of (19).

In nondivergence form, (12) becomes
$$(20) \qquad \rho_t + \rho v_x + \rho_x v = 0,$$
$$(21) \qquad \rho_t v + \rho v_t + \rho_x v^2 + 2\rho v v_x + p_x = 0.$$
Multiplication of (20) by $\sigma^2=p'(\rho)$ gives
$$(22) \qquad p_t + vp_x + \sigma^2\rho v_x = 0.$$
Equations (20)--(21) also give
$$(23) \qquad \rho v_t + \rho v v_x + p_x = 0.$$
Multiplying (23) by $\sigma$ and adding and subtracting it from (22) yields
$$(24) \qquad \begin{cases} p_t + (v+\sigma)p_x + \rho\sigma(v_t+(v+\sigma)v_x) = 0 \\ p_t + (v-\sigma)p_x - \rho\sigma(v_t+(v-\sigma)v_x) = 0. \end{cases}$$
Along (18)--(19), this becomes
$$(25) \qquad \begin{cases} \frac{d}{dt}[p(x_1(t),t)] + \rho(x_1(t),t)\sigma(x_1(t),t)\frac{d}{dt}[v(x_1(t),t)] = 0 \\ \frac{d}{dt}[p(x_2(t),t)] - \rho(x_2(t),t)\sigma(x_2(t),t)\frac{d}{dt}[v(x_2(t),t)] = 0. \end{cases}$$
Since $\frac{dp}{dt}=\sigma^2\frac{d\rho}{dt}$,
$$(26) \qquad \frac{\sigma}{\rho}\frac{d\rho}{dt} \mathbin{\pm} \frac{dv}{dt} = 0 \quad \text{along the trajectories of (18), (19),}$$
provided $\rho>0$.

For a Riemann invariant $w^1=w^1(\rho,v)$ along $x_1(\cdot)$,
$$0 = \frac{d}{dt}[w^1(\rho(x_1(t),t),v(x_1(t),t))] = \frac{\partial w^1}{\partial\rho}\frac{d}{dt}[\rho(x_1(t),t)] + \frac{\partial w^1}{\partial v}\frac{d}{dt}[v(x_1(t),t)].$$
Equation (26) gives
$$\frac{\partial w^1}{\partial\rho} = \frac{\sigma(\rho)}{\rho}, \quad \frac{\partial w^1}{\partial v} = 1.$$
and similarly
$$\frac{\partial w^2}{\partial\rho} = \frac{\sigma(\rho)}{\rho}, \quad \frac{\partial w^2}{\partial v} = -1.$$
Thus, up to additive constants, the Riemann invariants are
$$w^1 = \int_1^\rho \frac{\sigma(s)}{s}\,ds + v, \quad w^2 = \int_1^\rho \frac{\sigma(s)}{s}\,ds - v.$$
\end{example}

\subsection{Nonexistence of smooth solutions}\label{sec:nonexistence-smooth-solutions-systems}

\begin{theorem}[Riemann invariants and blow-up]
\label{thm:riemann-invariants-blowup}
\uses{def:riemann-invariant, thm:conservation-laws-riemann-invariants}
\dcref{ch11:11.3-thm-2}
\group{Conservation Laws}

Assume $g$ is smooth, with compact support. Suppose also the genuine nonlinearity condition
$$(27) \qquad \frac{\partial\lambda_i}{\partial w_j} > 0 \text{ in } \mathbb{R}^2 \quad (i=1,2,\ i\neq j)$$
holds. Then the initial-value problem (1) cannot have a smooth solution $u$ existing for all times $t\geq 0$ if
$$(28) \qquad \text{either } v^1_x < 0 \text{ or } v^2_x < 0 \text{ somewhere on } \mathbb{R}\times\{t=0\}.$$
\end{theorem}

\begin{proof}
\sketch Set $a=v^1_x$ and differentiate the first equation of (6). Along a $\lambda_2$-characteristic, introduce the integrating factor
\[
\xi(t)=\exp\!\left(\int_0^t
\frac{\partial_{w_2}\lambda_2}{\lambda_2-\lambda_1}
(v^2_t+\lambda_2v^2_x)\,ds\right).
\]
Because $v^1$ is constant on this characteristic, $\xi$ is bounded above and away from zero while $v$ remains bounded. The differentiated equation integrates to
\[
(35)\qquad
a(t)=a(0)\xi(t)^{-1}
\left(1+a(0)\int_0^t\partial_{w_1}\lambda_2\,\xi(s)^{-1}\,ds\right)^{-1}.
\]
Under (27), a negative $a(0)$ forces the denominator to vanish in finite time. Interchanging the two families gives the same conclusion when $v^2_x<0$ initially.
\end{proof}

\section{Entropy Criteria}\label{sec:entropy-criteria-systems}

The viscous approximation is
\[
(2)\qquad
\mathbf u_t^\varepsilon+\mathbf F(\mathbf u^\varepsilon)_x
-\varepsilon\mathbf u_{xx}^\varepsilon=0
\quad\text{in }\mathbb R\times(0,\infty).
\]

\subsection{Vanishing viscosity, traveling waves}\label{sec:vanishing-viscosity-traveling-waves}

A traveling profile has the form
\[
(3)\qquad \mathbf u^\varepsilon(x,t)
=\mathbf v\!\left(\frac{x-\sigma t}{\varepsilon}\right),
\]
with
\[
(5)\qquad
\lim_{s\to-\infty}\mathbf v(s)=u_l,
\quad \lim_{s\to+\infty}\mathbf v(s)=u_r,
\quad \lim_{s\to\pm\infty}\mathbf v'(s)=0.
\]
Substitution and one integration give
\[
(10)\qquad
\mathbf v'=\mathbf F(\mathbf v)-\mathbf F(u_l)-\sigma(\mathbf v-u_l),
\]
and the end states necessarily satisfy
\[
(9)\qquad \mathbf F(u_l)-\mathbf F(u_r)=\sigma(u_l-u_r).
\]

\begin{theorem}[Existence of traveling waves for genuinely nonlinear systems]
\label{thm:existence-traveling-waves-genuinely-nonlinear}
\uses{def:genuinely-nonlinear-linearly-degenerate, def:shock-set, thm:structure-of-shock-set}
\dcref{ch11:11.4-thm-1}
\group{Conservation Laws}

Assume the pair $(\lambda_k, \mathbf{r}_k)$ is genuinely nonlinear for $k=1,\ldots,m$. Let $u_r$ be selected sufficiently close to $u_l$. Then there exists a traveling wave solution of (2) connecting $u_l$ to $u_r$ if and only if
$$(12) \qquad u_r \in S_k(u_l)$$
for some $k\in\{1,\ldots,m\}$.
\end{theorem}

\begin{proof}
\sketch Any profile satisfying (5) obeys the Rankine--Hugoniot identity (9), hence $u_r\in S_k(u_l)$ for some $k$. Set
$\mathbf G(z)=\mathbf F(z)-\mathbf F(u_l)-\sigma(z-u_l)$, so the profile equation is $\mathbf v'=\mathbf G(\mathbf v)$ and the eigenvalues of $D\mathbf G(u_l)$ are $\lambda_j(u_l)-\sigma$. By \cref{thm:structure-of-shock-set}(iii),
\[
\lambda_k(u_l)-\sigma
=\frac{\lambda_k(u_l)-\lambda_k(u_r)}{2}+o(|u_r-u_l|).
\]
An orbit approaching $u_l$ as $s\to-\infty$ requires $\lambda_k(u_l)-\sigma>0$. For nearby states this gives $\lambda_k(u_r)<\lambda_k(u_l)$, hence $u_r\in S_k^-(u_l)$. The source omits the sufficiency of (12) and refers its proof to \cite{MajdaPego1985}.
\end{proof}

\begin{definition}[Lax's entropy criterion, general form]
\label{def:lax-entropy-criterion-general}
\uses{def:shock-set}
\dcref{ch11:11.4-def-1}
\group{Conservation Laws}

$$(16) \qquad \begin{cases} \sigma(z,u_l) > \sigma(u_r,u_l) \text{ for each } z \text{ lying} \\ \text{on the curve } S_k(u_l) \text{ between } u_r \text{ and } u_l. \end{cases}$$
Condition (16) is \textit{Lax's entropy criterion}. It holds automatically for sufficiently close states on a genuinely nonlinear admissible shock branch.
\end{definition}

\begin{example}[Traveling waves for the $p$-system]
\label{ex:traveling-waves-p-system}
\uses{ex:p-system-definition, thm:existence-traveling-waves-genuinely-nonlinear, def:lax-entropy-criterion-general}
\dcref{ch11:11.4-ex-1}
\group{Conservation Laws}

The $p$-system is
$$(17) \qquad \begin{cases} u^1_t - u^2_x = 0 & \text{(compatibility condition)} \\ u^2_t - p(u^1)_x = 0 & \text{(Newton's law)}, \end{cases}$$
under the usual strict hyperbolicity condition
$$(18) \qquad p' > 0.$$
Its regularization is
$$(19) \qquad \begin{cases} u^{\varepsilon,1}_t - u^{\varepsilon,2}_x = 0 \\ u^{\varepsilon,2}_t - p(u^{\varepsilon,1})_x = \varepsilon u^{\varepsilon,2}_{xx}. \end{cases}$$
Let $u^\varepsilon = v\left(\frac{x-\sigma t}{\varepsilon}\right)$ satisfy
$$(20) \qquad \lim_{s\to-\infty} v = u_l, \quad \lim_{s\to\infty} v = u_r, \quad \lim_{s\to\pm\infty} \dot v = 0.$$
For $\mathbf v=(v^1,v^2)$, equation (19) becomes
$$(21) \qquad \begin{cases} -\sigma\dot v^1 - \dot v^2 = 0 & \left(\cdot = \frac{d}{ds}\right) \\ -\sigma\dot v^2 - p(v^1)^\cdot = \ddot v^2. \end{cases}$$
Integration using (20) gives
$$(22) \qquad \begin{cases} \sigma v^1 + v^2 = \sigma v^1_l + v^2_l = \sigma v^1_r + v^2_r \\ \dot v^2 = \sigma(v^2_l-v^2) + p(v^1_l) - p(v^1) = \sigma(v^2_r-v^2) + p(v^1_r) - p(v^1), \end{cases}$$
for $u_l=(v^1_l,v^2_l)$, $u_r=(v^1_r,v^2_r)$. In particular,
$$\begin{cases} \sigma v^1_l + v^2_l = \sigma v^1_r + v^2_r \\ \sigma v^2_l + p(v^1_l) = \sigma v^2_r + p(v^1_r). \end{cases}$$
Thus
$$(23) \qquad \sigma^2 = \frac{p(v^1_r)-p(v^1_l)}{v^1_r - v^1_l}.$$
If $v^1_r>v^1_l$ and $\sigma>0$, the entropy criterion is
$$(24) \qquad \frac{p(z_1)-p(v^1_l)}{z_1-v^1_l} > \frac{p(v^1_r)-p(v^1_l)}{v^1_r-v^1_l}$$
for all $z$ on the curve $S_k(u_l)$ between $u_l$ and $u_r$, $z=(z_1,z_2)$.

Eliminating $v^2$ from (22) gives
$$\dot v^1 = \frac{p(v^1)-p(v^1_l)}{\sigma} - \sigma(v^1-v^1_l) =: g(v^1).$$
By (23), $g(v^1_l)=g(v^1_r)=0$. A connecting profile exists exactly when
$$g(z_1) > 0 \quad \text{for } v^1_l < z_1 < v^1_r.$$
which is (24); the case $v^1_r<v^1_l$ is analogous.
\end{example}

\subsection{Entropy/entropy-flux pairs}\label{sec:entropy-entropy-flux-pairs}

\begin{definition}[Entropy/entropy-flux pair, system case]
\label{def:entropy-entropy-flux-pair-system}
\dcref{ch11:11.4-def-2}
\group{Conservation Laws}
Two smooth functions $\Phi, \Psi:\mathbb{R}^m\to\mathbb{R}$ comprise an \textit{entropy/entropy-flux pair} for the conservation law $u_t + \mathbf{F}(u)_x = 0$ provided

$$(25) \qquad \Phi \text{ is convex}$$

and

$$(26) \qquad D\Phi(z)D\mathbf{F}(z) = D\Psi(z) \quad (z\in\mathbb{R}^m).$$
\end{definition}

For an entropy pair, the distributional entropy inequality is
\[
(28)\qquad \Phi(\mathbf u)_t+\Psi(\mathbf u)_x\leq0,
\]
meaning
\[
(29)\qquad
\int_0^\infty\!\int_{-\infty}^{\infty}
\bigl(\Phi(\mathbf u)v_t+\Psi(\mathbf u)v_x\bigr)\,dx\,dt\geq0
\]
for every nonnegative $v\in C_c^\infty(\mathbb R\times(0,\infty))$. The associated initial-value problem is
\[
(30)\qquad
\begin{cases}
\mathbf u_t+\mathbf F(\mathbf u)_x=0 & \text{in }\mathbb R\times(0,\infty),\\
\mathbf u=\mathbf g & \text{on }\mathbb R\times\{t=0\}.
\end{cases}
\]

\begin{definition}[Entropy solution, system case]
\label{def:entropy-solution-system}
\uses{def:integral-solution-system, def:entropy-entropy-flux-pair-system}
\dcref{ch11:11.4-def-3}
\group{Conservation Laws}

We call $u$ an \textit{entropy solution} of (30) provided $u$ is an integral solution and $u$ satisfies the inequalities (29) for each entropy/entropy-flux pair $(\Phi,\Psi)$.
\end{definition}

Let $\mathbf u^\varepsilon$ be smooth, rapidly decaying solutions of
\[
(31)\qquad
\begin{cases}
\mathbf u_t^\varepsilon+\mathbf F(\mathbf u^\varepsilon)_x
-\varepsilon\mathbf u_{xx}^\varepsilon=0,\\
\mathbf u^\varepsilon=\mathbf g & \text{on }\mathbb R\times\{t=0\},
\end{cases}
\]
uniformly bounded in $L^\infty$, and suppose
\[
(32)\qquad \mathbf u^\varepsilon\to\mathbf u\quad\text{a.e. as }\varepsilon\to0.
\]

\begin{theorem}[Entropy and vanishing viscosity]
\label{thm:entropy-vanishing-viscosity}
\uses{def:entropy-solution-system, def:entropy-entropy-flux-pair-system}
\dcref{ch11:11.4-thm-2}
\group{Conservation Laws}

The function $u$ is an entropy solution of the conservation law (30).
\end{theorem}

\begin{proof}
\sketch For any entropy pair, multiplication of (31) by $D\Phi(\mathbf u^\varepsilon)$ gives
\[
\Phi(\mathbf u^\varepsilon)_t+\Psi(\mathbf u^\varepsilon)_x
=\varepsilon\Phi(\mathbf u^\varepsilon)_{xx}
-\varepsilon\bigl(D^2\Phi(\mathbf u^\varepsilon)\mathbf u_x^\varepsilon\bigr)\cdot\mathbf u_x^\varepsilon.
\]
The last term is nonpositive by convexity. Testing against nonnegative $v$, integrating by parts, and using (32) with dominated convergence yields (29). Testing (31) against compactly supported vector fields and passing to the limit similarly gives the weak identity (13), so $\mathbf u$ is also an integral solution of (30).
\end{proof}

\begin{example}[Entropy pairs for a scalar conservation law]
\label{ex:scalar-entropy-flux-existence}
\uses{def:entropy-entropy-flux-pair-system}
\dcref{ch11:11.4-ex-2}
\group{Conservation Laws}

For $m=1$, any convex $\Phi$ has the corresponding flux
$$\Psi(z) = \int_{z_0}^z \Phi'(w)F'(w)\,dw \quad (z\in\mathbb{R}).$$
\end{example}

\begin{example}[Entropy pair for the $p$-system]
\label{ex:p-system-entropy-pair}
\uses{def:entropy-entropy-flux-pair-system, ex:p-system-definition}
\dcref{ch11:11.4-ex-3}
\group{Conservation Laws}

For the $p$-system we have $m=2$. To verify (25), (26) we must find $\Phi,\Psi$, with $\Phi$ convex and
$$(\Phi_{z_1},\Phi_{z_2})\begin{pmatrix} 0 & -1 \\ -p'(z_1) & 0 \end{pmatrix} = \begin{pmatrix} \Psi_{z_1} \\ \Psi_{z_2} \end{pmatrix}.$$
A solution is
$$\Phi(z) = \frac{z_2^2}{2} + \int_0^{z_1} p(w)\,dw, \quad \Psi(z) = -p(z_1)z_2 \quad (z\in\mathbb{R}^2).$$
The condition $p'>0$ makes $\Phi$ convex.
\end{example}

\subsection{Uniqueness for a scalar conservation law}\label{sec:uniqueness-scalar-conservation-law-entropy}

Let $F:\mathbb R\to\mathbb R$ be smooth and consider
\[
(35)\qquad
\begin{cases}
u_t+F(u)_x=0 & \text{in }\mathbb R\times(0,\infty),\\
u=g & \text{on }\mathbb R\times\{t=0\}.
\end{cases}
\]

\begin{definition}[Entropy/entropy-flux pair, scalar case]
\label{def:entropy-entropy-flux-pair-scalar}
\dcref{ch11:11.4-def-4}
\group{Conservation Laws}
Two smooth functions $\Phi,\Psi:\mathbb{R}\to\mathbb{R}$ comprise an \textit{entropy/entropy-flux pair} for the conservation law $u_t + F(u)_x = 0$ provided

$$(36) \qquad \Phi \text{ is convex}$$

and

$$(37) \qquad \Phi'(z)F'(z) = \Psi'(z) \quad (z\in\mathbb{R}).$$
\end{definition}

For a scalar entropy pair, the entropy inequality means
\[
(38)\qquad
\int_0^\infty\!\int_{-\infty}^{\infty}
\bigl(\Phi(u)v_t+\Psi(u)v_x\bigr)\,dx\,dt\geq0
\]
for every nonnegative $v\in C_c^\infty(\mathbb R\times(0,\infty))$.

\begin{definition}[Entropy solution, scalar case]
\label{def:entropy-solution-scalar-new}
\uses{def:entropy-entropy-flux-pair-scalar}
\dcref{ch11:11.4-def-5}
\group{Conservation Laws}

We call $u\in C([0,\infty), L^1(\mathbb{R}))\cap L^\infty(\mathbb{R}\times(0,\infty))$ an \textit{entropy solution} of (35) provided $u$ satisfies the inequalities (38) for each entropy/entropy-flux pair $(\Phi,\Psi)$, and $u(\cdot,t)\to g$ in $L^1$ as $t\to 0$.
\end{definition}

\begin{remark}[Consistency with the scalar entropy solution of Chapter 3]
\label{rem:entropy-solution-scalar-consistency}
\uses{def:entropy-solution-scalar-new, def:integral-solution-conservation-law}
\dcref{ch11:11.4-rem-1}
\group{Conservation Laws}

Taking $\Phi(z)=\pm z$ and $\Psi(z)=\pm F(z)$ in (38) gives
$$\int_0^\infty\int_{-\infty}^\infty uv_t + F(u)v_x\,dx\,dt = 0$$
for all $v\in C_c^1(\mathbb{R}\times(0,\infty))$. Since $u(\cdot,t)\to g$ in $L^1$,
$$\int_0^\infty\int_{-\infty}^\infty uv_t + F(u)v_x\,dx\,dt + \int_{-\infty}^\infty gv\,dx\Big|_{t=0} = 0$$
for all $v\in C_c^1(\mathbb{R}\times[0,\infty))$. Thus an entropy solution in this sense is an integral solution, consistently with \cref{def:integral-solution-conservation-law}.
\end{remark}

\begin{theorem}[Uniqueness of entropy solutions for a single conservation law]
\label{thm:uniqueness-entropy-solutions-scalar-new}
\uses{def:entropy-solution-scalar-new, thm:uniqueness-viscosity-solution-hj}
\dcref{ch11:11.4-thm-3}
\group{Conservation Laws}

There exists---up to a set of measure zero---at most one entropy solution of (35).
\end{theorem}

\begin{proof}
\sketch Approximate $|z-\alpha|$ by smooth convex entropies. Passing to the limit in (38) gives, for every $\alpha\in\mathbb R$, the Kruzkov inequality with entropy flux
$\operatorname{sgn}(u-\alpha)(F(u)-F(\alpha))$. Apply this inequality to $u$ with $\alpha=\bar u(y,s)$ and to another entropy solution $\bar u$ with $\alpha=u(x,t)$, then add and integrate in the doubled variables $(x,t,y,s)$. Testing with mollifiers concentrating on $x=y$ and $t=s$ yields the distributional inequality
\[
\partial_t|u-\bar u|
+\partial_x\!\left(\operatorname{sgn}(u-\bar u)(F(u)-F(\bar u))\right)\leq0.
\]
Spatial cutoffs tending to one and time cutoffs at $s<t$ imply
\[
\|u(\cdot,t)-\bar u(\cdot,t)\|_{L^1}
\leq\|u(\cdot,s)-\bar u(\cdot,s)\|_{L^1}
\]
for almost every $s<t$. Letting $s\downarrow0$ and using the common initial trace $g$ proves equality almost everywhere.
\end{proof}

\section{References}\label{sec:chapter11-references}
\textbf{Hyperbolic and shock systems.} \cite{Zheng2001}, \cite{Conlon1980}, \cite{LeVeque1992}, \cite{CourantFriedrichs1976}, \cite{XiaoZhang1989}, \cite{Zheng2001}.
\textbf{Riemann invariants.} \cite{Logan1987}.
\textbf{Viscous waves and entropy.} \cite{MajdaPego1985}, \cite{Conlon1980}, \cite{Smoller1983}.
\textbf{Further problems.} \cite{Kato1980}.
